\documentclass[11pt]{article}
\usepackage{amsmath,amssymb,amsthm}
\usepackage[margin=1in]{geometry}

\newtheorem{theorem}{Theorem}
\newtheorem{lemma}[theorem]{Lemma}

\title{Problem 439: Sum of Sum of Divisors}
\author{Project Euler Solutions}
\date{}

\begin{document}
\maketitle

\section{Problem Statement}
Let $\sigma(k)$ denote the sum of divisors of $k$. Define
\[
S(N) = \sum_{i=1}^{N}\sum_{j=1}^{N} \sigma(i \cdot j).
\]
Given $S(10^3) = 563576517282$ and $S(10^5) \bmod 10^9 = 215766508$, find $S(10^{11}) \bmod 10^9$.

\section{Mathematical Foundation}

\begin{theorem}[Multiplicative Convolution Identity]
For all positive integers $m, n$:
\[
\sigma(mn) = \sum_{d \mid \gcd(m,n)} \mu(d) \cdot \sigma(m/d) \cdot \sigma(n/d).
\]
\end{theorem}

\begin{proof}
Both sides are multiplicative in the pair $(m, n)$, so it suffices to verify for prime powers. Let $m = p^a$, $n = p^b$ with $a \leq b$. Then:
\[
\text{RHS} = \sum_{k=0}^{a} \mu(p^k)\,\sigma(p^{a-k})\,\sigma(p^{b-k}) = \sigma(p^a)\sigma(p^b) - \sigma(p^{a-1})\sigma(p^{b-1}).
\]
Using $\sigma(p^c) = (p^{c+1}-1)/(p-1)$, direct computation confirms this equals $\sigma(p^{a+b})$.
\end{proof}

\begin{theorem}[Main Summation Formula]
Define $T(n) = \sum_{k=1}^n \sigma(k)$. Then
\[
S(N) = \sum_{d=1}^{N} \mu(d) \cdot T\!\left(\left\lfloor \frac{N}{d}\right\rfloor\right)^{\!2}.
\]
\end{theorem}

\begin{proof}
Substitute Theorem~1 into $S(N)$ and set $i = da$, $j = db$:
\[
S(N) = \sum_{d=1}^N \mu(d) \sum_{a=1}^{\lfloor N/d \rfloor} \sigma(a) \sum_{b=1}^{\lfloor N/d \rfloor} \sigma(b) = \sum_{d=1}^N \mu(d)\, T(\lfloor N/d \rfloor)^2. \qedhere
\]
\end{proof}

\begin{lemma}[Efficient $T(n)$]
\[
T(n) = \sum_{d=1}^n d \cdot \left\lfloor \frac{n}{d}\right\rfloor,
\]
computable in $O(\sqrt{n})$ by grouping equal values of $\lfloor n/d \rfloor$.
\end{lemma}

\begin{proof}
$T(n) = \sum_{k=1}^n \sum_{d \mid k} d = \sum_{d=1}^n d \lfloor n/d \rfloor$ by swapping summation. The function $\lfloor n/d \rfloor$ takes $O(\sqrt{n})$ distinct values, and for each block, $\sum d$ is a closed-form arithmetic sum.
\end{proof}

\begin{theorem}[Sub-Linear Mertens Function]
The Mertens function $M(n) = \sum_{k=1}^n \mu(k)$ satisfies
\[
M(n) = 1 - \sum_{k=2}^{n} M\!\left(\left\lfloor \frac{n}{k}\right\rfloor\right)
\]
and is computable in $O(n^{2/3})$ time and space.
\end{theorem}

\begin{proof}
From $\sum_{d=1}^n \mu(d)\lfloor n/d \rfloor = 1$ (a consequence of $\sum_{d \mid k} \mu(d) = [k=1]$), rearranging gives the recursion. Sieving $\mu$ up to $n^{2/3}$ and recursing on $O(n^{1/3})$ large arguments yields $O(n^{2/3})$ total complexity.
\end{proof}

\section{Algorithm}
\begin{enumerate}
    \item Sieve $\mu(d)$ for $d \leq N^{2/3}$ and compute prefix sums.
    \item Implement $M(n)$ using the sieve for small $n$ and the recursion with memoization for large $n$.
    \item Implement $T(n) \bmod p$ via the $O(\sqrt{n})$ block summation.
    \item Evaluate $S(N) = \sum_d \mu(d)\,T(\lfloor N/d \rfloor)^2$ by grouping $d$-values with equal $\lfloor N/d \rfloor$, using $M$ for partial sums of $\mu$.
\end{enumerate}

\section{Complexity Analysis}
\begin{itemize}
    \item \textbf{Time:} $O(N^{2/3})$ dominated by the Mertens sieve and recursion.
    \item \textbf{Space:} $O(N^{2/3})$ for the sieve and memoization table.
\end{itemize}

\section{Answer}

\[
\boxed{968697378}
\]

\end{document}
